\documentclass[10pt,a4paper]{amsart}
\usepackage[margin=19mm]{geometry}
\usepackage{amsmath,amssymb,mathtools}
\usepackage[scr=rsfs,cal=euler]{mathalfa}
\usepackage{xfrac}
\usepackage[unicode,hidelinks,bookmarks=false]{hyperref}
\newcommand{\ie}{i.e.\ }
\newcommand{\cf}{cf.\ }
\newcommand{\resp}{respectively\ }
\newcommand{\Subsection}{\subsection}
\providecommand{\nobreakdash}{\nobreak\hbox{-}}
\setlength{\emergencystretch}{2em}
\multlinegap0pt
\usepackage{amssymb}
\usepackage{enumerate}
\usepackage{tikz}
\usepackage{float}
\usepackage{mathtools}
\newtheorem{theorem}{Theorem}
\title{On solution of Fractional Diffusion Equation using Conformable Laplace Transform}
\author{Somnath Sarate, Anil Khairnar, Krishnat Masalkar}
\date{Source: arXiv:2605.12157v1 (CC BY 4.0)}
\begin{document}
\maketitle

\noindent\emph{Source and licence.} On solution of Fractional Diffusion Equation using Conformable Laplace Transform. Version arXiv:2605.12157v1 (CC BY 4.0), distributed by arXiv under the Creative Commons Attribution 4.0 International licence, http://creativecommons.org/licenses/by/4.0/, as recorded in the arXiv licence metadata for this exact version and shown on the abstract page https://arxiv.org/abs/2605.12157v1. Full licence notice: \texttt{licenses/arxiv-2605.12157-CC-BY-4.0.txt}.\par\smallskip

\noindent\textit{Editorial scope.} Counted unit: the article's theorem theorem (source0.tex lines 402--416). The article's complete original proof of it is delivered with it (source0.tex lines 417--535, one \texttt{proof} environment of 119 lines). No mathematics is written, repaired or re-derived: the statement and the proof are the article's own lines, in the article's own notation, and no retained line is substituted or re-flowed.  Retained uncounted as the source-local context the proof needs: lines 112--117.  Nothing was omitted from any retained span, so the package carries no omission record.  No retained line contains a citation, so the wrapper carries no bibliography and no bibliography entry.  No retained line contains a figure environment, an image-inclusion command or drawing code, so no image is delivered and the package declares no asset dependency.  Editorial formatting only, in addition: an AMS \texttt{amsart} preamble that restates the article's own declarations and macros used by the retained lines, namely the environments \texttt{theorem} on separate counters, together with the standard packages those lines need.  The retained environments print in this document as Theorem 1 in order of appearance, whereas the article prints the same items with the numbering of its own full document.  The complete native source of the article as distributed in the arXiv e-print for this exact version is bundled unchanged as \texttt{sources/200404/source0.tex}.
\begin{theorem}
\label{thm:2.1}
	Let $a\in \mathbb R$ and $\phi_a:[a,\infty)\rightarrow \mathbb R$ be the function defined by
	$\phi_a(t)=\frac{(t-a)^{\alpha}}{\alpha}.$
	Let $f:[a, \infty) \rightarrow \mathbb C$ be a function. Then $\mathcal{L}_{a}^{\alpha}(f(t))(s)$ exist if and only if $\mathcal{L}((f\circ\phi_a^{-1})(t))(s)$ exist on $[a,\infty)$ for all $s$ with $real(s)>b$ for some $b\in \mathbb R$. 
\end{theorem}

\begin{theorem}
Let $f(t)$ be $n$-times differentiable on $[0,\infty)$ and let $0<\alpha,\beta\leq 1$. Then the conformable fractional Laplace transform of the $n$-th conformable derivative of order $\beta$ is
\begin{align*}
\mathcal{L}_0^{\alpha} \left\{ (T_\beta)^n f(t) \right\} = s^n \mathcal{L}_0^{\alpha} \left\{ t^{n(\alpha-\beta)} f(t) \right\} 
&+ \sum_{k=1}^n \binom{n}{k} (\beta-\alpha)^k s^{\,n-k} \mathcal{L}_0^{\alpha} \left\{ t^{(n-k)(\alpha-\beta)-k\beta} f(t) \right\} \\
+ \text{boundary terms}.
\end{align*}
where the boundary terms are
$
f(0), T_\beta f(0), \ldots, (T_\beta)^{\,n-1} f(0),
$ and the conformable Laplace transform is
\[
\mathcal{L}_0^{\alpha} \{ f(t) \} = \int_0^\infty e^{-s \frac{t^\alpha}{\alpha}} t^{\alpha-1} f(t) \, dt.
\]
\end{theorem}

\begin{proof}
    From Theorem~\ref{thm:2.1}, for the first conformable derivative, we have
    \[
\mathcal{L}_{o}^{\alpha}\{T_{\beta}f(t)\}
= \left[e^{-s\frac{t^{\alpha}}{\alpha}}t^{\alpha-\beta}f(t)\right]_{0}^{\infty}+
s\mathcal{L}_{o}^{\alpha}\{t^{\alpha-\beta}f(t)\}
+
(\beta-\alpha)\mathcal{L}_{o}^{\alpha}\{t^{-\beta}f(t)\}
\]  
Applying the conformable fractional Laplace transform again,
\[
\mathcal{L}_0^{\alpha}\{(T_{\beta})^{2}f(t)\}
=
\mathcal{L}_0^{\alpha}\{T_{\beta}(T_{\beta}f(t))\}.
\]
Using the first derivative formula with
\(
g(t)=T_{\beta}f(t),
\)
\[
\mathcal{L}_0^{\alpha}\{T_{\beta}g(t)\}
= \left[e^{-s\frac{t^{\alpha}}{\alpha}}t^{\alpha-\beta}g(t)\right]_{0}^{\infty}+
s\mathcal{L}_0^{\alpha}\{t^{\alpha-\beta}g(t)\}
+
(\beta-\alpha)\mathcal{L}_0^{\alpha}\{t^{-\beta}g(t)\}.
\]
Hence,
\[
\mathcal{L}_0^{\alpha}\{(T_{\beta})^{2}f(t)\}
= \left[e^{-s\frac{t^{\alpha}}{\alpha}}t^{\alpha-\beta}T_{\beta}f(t)\}\right]_{0}^{\infty}+
s\mathcal{L}_0^{\alpha}\{t^{\alpha-\beta}T_{\beta}f(t)\}
+
(\beta-\alpha)\mathcal{L}_0^{\alpha}\{t^{-\beta}T_{\beta}f(t)\}.
\]
Using the definition of the conformable derivative
\(
T_{\beta}f(t)=t^{1-\beta}f'(t),
\)
we obtain
\[
t^{\alpha-\beta}T_{\beta}f(t)
=t^{\alpha-\beta}\,t^{1-\beta}f'(t)=t^{\alpha+1-2\beta}f'(t),
\quad\text{and}\quad t^{-\beta}T_{\beta}f(t) = t^{1-2\beta}f'(t).
\]
Thus,
\begin{align*}
\mathcal{L}_0^{\alpha}\{(T_{\beta})^{2}f(t)\}
&= \left[e^{-s\frac{t^{\alpha}}{\alpha}}t^{\alpha-\beta}t^{\alpha-\beta}T_{\beta}f(t)\}\right]_{0}^{\infty}\\
& +
s\mathcal{L}_0^{\alpha}\{t^{\alpha+1-2\beta}f'(t)\}\\
&+
(\beta-\alpha)\mathcal{L}_0^{\alpha}\{t^{1-2\beta}f'(t)\}.
\end{align*}
Applying integration by parts to the above transforms, we obtain
\[
\mathcal{L}_0^{\alpha}\{t^{\alpha+1-2\beta}f'(t)\}
=
\left[e^{-\frac{s t^{\alpha}}{\alpha}} t^{2\alpha-2\beta}f(t)\right]_{0}^{\infty}
+s\mathcal{L}_0^{\alpha}\{t^{2(\alpha-\beta)}f(t)\}
+(\beta-\alpha)\mathcal{L}_0^{\alpha}\{t^{\alpha-2\beta}f(t)\},
\]
and
\[
\mathcal{L}_0^{\alpha}\{t^{1-2\beta}f'(t)\}
=
\left[e^{-\frac{s t^{\alpha}}{\alpha}} t^{\alpha-2\beta}f(t)\right]_{0}^{\infty}
+s\mathcal{L}_0^{\alpha}\{t^{\alpha-2\beta}f(t)\}
+(\beta-\alpha)\mathcal{L}_0^{\alpha}\{t^{-2\beta}f(t)\}.
\]
Substituting these expressions back, we obtain
\begin{align*}
\mathcal{L}_0^{\alpha}\{(T_{\beta})^{2}f(t)\}
&= \left[e^{-s\frac{t^{\alpha}}{\alpha}}t^{\alpha-\beta}t^{\alpha-\beta}T_{\beta}f(t)\}\right]_{0}^{\infty}\\
& + s\left[e^{-\frac{s t^{\alpha}}{\alpha}} t^{2\alpha-2\beta}f(t)\right]_{0}^{\infty}
+s^2\mathcal{L}_0^{\alpha}\{t^{2(\alpha-\beta)}f(t)\}
+s(\beta-\alpha)\mathcal{L}_0^{\alpha}\{t^{\alpha-2\beta}f(t)\}\\
&+ (\beta-\alpha)\left[e^{-\frac{s t^{\alpha}}{\alpha}} t^{\alpha-2\beta}f(t)\right]_{0}^{\infty}
+s(\beta-\alpha)\mathcal{L}_0^{\alpha}\{t^{\alpha-2\beta}f(t)\}\\
&+(\beta-\alpha)^2\mathcal{L}_0^{\alpha}\{t^{-2\beta}f(t)\}.
\end{align*}
\begin{align*}
\mathcal{L}_0^{\alpha}\{(T_{\beta})^{2}f(t)\}
&=s^2\mathcal{L}_0^{\alpha}\{t^{2(\alpha-\beta)}f(t)\}
+s(\beta-\alpha)\mathcal{L}_0^{\alpha}\{t^{\alpha-2\beta}f(t)
+s(\beta-\alpha)\mathcal{L}_0^{\alpha}\{t^{\alpha-2\beta}f(t)\}\\
&+(\beta-\alpha)^2\mathcal{L}_0^{\alpha}\{t^{-2\beta}f(t)\}
+ \text{boundary terms}
\end{align*}
where boundary terms are
\begin{align*}
 (\beta-\alpha)\left[e^{-\frac{s t^{\alpha}}{\alpha}} t^{\alpha-2\beta}f(t)\right]_{0}^{\infty},\quad
 s\left[e^{-\frac{s t^{\alpha}}{\alpha}} t^{2\alpha-2\beta}f(t)\right]_{0}^{\infty},
 \quad\left[e^{-s\frac{t^{\alpha}}{\alpha}}t^{\alpha-\beta}t^{\alpha-\beta}T_{\beta}f(t)\}\right]_{0}^{\infty}
\end{align*}
Since the exponential term vanishes as $t \to \infty$, the boundary contributions reduce to the initial values involving
$
f(0)$, $T_{\beta}f(0).$
 Therefore, the conformable fractional Laplace transform of the second conformable derivative is given by
\[
\mathcal{L}_0^{\alpha}\{(T_{\beta})^{2}f(t)\}
=
s^{2}\mathcal{L}_0^{\alpha}\{t^{2(\alpha-\beta)}f(t)\}
+2(\beta-\alpha)s\mathcal{L}_0^{\alpha}\{t^{\alpha-2\beta}f(t)\}
+(\beta-\alpha)^{2}\mathcal{L}_0^{\alpha}\{t^{-2\beta}f(t)\},
\]
together with boundary terms involving $f(0)$ and $T_{\beta}f(0)$ where $0<\alpha,\beta\le1$.\\
Continuing this procedure, each application of the conformable fractional Laplace transform introduces an additional factor of $s$ and additional terms involving powers of $(\beta-\alpha)$. The coefficients follow the binomial pattern arising from repeated application of the transform.
Thus, after $n$ iterations we obtain
\begin{align*}
\mathcal{L}_0^{\alpha} \left\{ (T_\beta)^n f(t) \right\} = s^n \mathcal{L}_0^{\alpha} \left\{ t^{n(\alpha-\beta)} f(t) \right\} 
&+ \sum_{k=1}^n \binom{n}{k} (\beta-\alpha)^k s^{\,n-k} \mathcal{L}_0^{\alpha} \left\{ t^{(n-k)(\alpha-\beta)-k\beta} f(t) \right\} \\
+ \text{boundary terms}.
\end{align*}
where the boundary terms are
$
f(0), T_\beta f(0), \ldots, (T_\beta)^{\,n-1} f(0).
$\\
Hence proved. 
\end{proof}

\end{document}
